\documentclass[aspectratio=169,10pt]{beamer}

\usetheme{Madrid}
\usecolortheme{seahorse}
\setbeamertemplate{navigation symbols}{}
\setbeamertemplate{footline}[frame number]

\usepackage{amsmath,amssymb,booktabs,array}
\usepackage{tikz}
\usetikzlibrary{arrows.meta,positioning}
\usepackage{xcolor}
\usepackage{hyperref}

\definecolor{darkblue}{RGB}{0,65,120}
\hypersetup{colorlinks=true,linkcolor=darkblue,urlcolor=darkblue}

\newcommand{\code}[1]{\texttt{#1}}
\newcommand{\TODO}[1]{\textcolor{red}{\textbf{TODO:}} \emph{#1}}

\title[L03: Balanced PF]{L03 -- Balanced (Single-phase Equivalent) Power Flow}
\subtitle{Foundations Lecture Series}
\author{Course Instructor}
\institute{Microgrid 3-phase Security AI Project}
\date{Foundations Lecture 3}

\begin{document}

\begin{frame}
  \titlepage
  \begin{center}
    \small Keywords: positive sequence, single-line diagram, three-phase symmetry, when balance holds
  \end{center}
\end{frame}

% ============================================================
\begin{frame}{Lecture roadmap}
\begin{enumerate}
  \item What ``balanced'' actually means
  \item The single-phase equivalent: how we collapse three phases into one
  \item Positive-sequence networks
  \item When the balanced assumption is appropriate
  \item When it lies: preview of L04
\end{enumerate}
\vspace{0.6em}
\begin{block}{Prerequisite}
L01 -- you should be comfortable with $Y_{\text{bus}}$ and the PF equations. L02 helpful but not required.
\end{block}
\end{frame}

% ============================================================
\section*{1. What ``balanced'' means}

\begin{frame}{Definition of balanced operation}
\TODO{Three-phase voltages and currents equal in magnitude, $120^\circ$ apart, sinusoidal. Source is symmetric, loads are symmetric, line parameters are symmetric. Visual: phasor diagram of a balanced set.}
\end{frame}

\begin{frame}{Symmetry assumptions in three places}
\TODO{(i) Source symmetry, (ii) load symmetry, (iii) line geometry symmetry (or transposition). All three usually hold in transmission, often fail in distribution. Bridge: under all three, three coupled equations collapse to one.}
\end{frame}

% ============================================================
\section*{2. The single-phase equivalent}

\begin{frame}{Collapsing three phases into one}
\TODO{Show that under balance, $V_b = a^2 V_a$ and $V_c = a V_a$ with $a = e^{j 2\pi/3}$. Therefore we only need to solve for phase $a$ -- the other phases follow by rotation. This is the ``single-line diagram'' you usually see.}
\end{frame}

\begin{frame}{Single-line diagrams in practice}
\TODO{Visual example: a small system shown both as a true three-phase circuit and as its single-line equivalent. Note that one-line diagrams in textbooks and pandapower's \code{runpp} both implicitly use this collapse.}
\end{frame}

% ============================================================
\section*{3. Positive-sequence networks}

\begin{frame}{Positive sequence in 30 seconds}
\TODO{Forward-reference to L04's symmetrical components: positive, negative, zero. Under balance, only positive sequence is non-zero. Negative and zero sequence appear only when balance breaks.}
\end{frame}

\begin{frame}{The math is L01's math}
\TODO{Show that the PF equations from L01 \emph{are} the positive-sequence PF equations. \code{pp.runpp} solves them. Nothing new computationally -- only the interpretation changes.}
\end{frame}

% ============================================================
\section*{4. When balance is appropriate}

\begin{frame}{When the assumption is fine}
\TODO{Transmission grids: three-phase loads, transposed lines, large industrial loads. The assumption introduces small ($< 1\%$) errors. List several practical examples where balanced PF is the right tool.}
\end{frame}

\begin{frame}{Quick-and-dirty hosting analyses}
\TODO{For distribution-level studies where the question is ``rough voltage profile, rough loading'', balanced PF is sometimes acceptable as a first pass. Note caveats: never claim phase-specific results from a balanced model.}
\end{frame}

% ============================================================
\section*{5. When it lies}

\begin{frame}{Where the balanced assumption breaks}
\TODO{Single-phase loads, single-phase rooftop PV, EV charging, asymmetric line geometry, fault analysis. Show a small example where balanced PF reports ``all good'' but the actual three-phase analysis (L04) shows phase A undervoltage.}
\end{frame}

\begin{frame}{Diagnostic: when do I need L04?}
\TODO{A short checklist for students: (1) Are loads single-phase? (2) Is PV single-phase? (3) Are line parameters asymmetric? (4) Is the network LV? If any yes, you need unbalanced PF.}
\end{frame}

% ============================================================
\begin{frame}{Homework}
\TODO{Take the L01 three-bus system, intentionally make one of the loads single-phase, and explain why \code{runpp} silently lies. Run \code{runpp\_3ph} on the same system and report the per-phase voltage spread.}
\end{frame}

\begin{frame}{Exit checklist}
\TODO{Items: (1) I can state the three symmetry assumptions, (2) I can draw a one-line from a three-phase, (3) I can decide whether to use \code{runpp} or \code{runpp\_3ph}, (4) I know the connection to symmetrical components is coming in L04.}
\end{frame}

\begin{frame}{References}
\TODO{Grainger \& Stevenson chapter on per-phase analysis, IEC/IEEE definitions of balanced operation, pandapower ``AC Power Flow'' page.}
\end{frame}

\end{document}
